\documentclass{article}
\usepackage{tcee-repaso}
\tcetema[Teoría de juegos y subastas]{3.A.14}{Teoría de juegos. Principales conceptos. Aplicaciones, con especial referencia a las subastas.}
\tcefecha{21/10/2026}
\begin{document}
\cabecerarepaso
\begin{columnas}

% ==========================================================================
\bloque{Esquema del cante}
\esquemacante{%
\apartadocante{Introducción}{2,5}
\apartadocante{I. Conceptos (I): información completa}{9}
\subcante{1. Elementos, representación y clasificación}{1,5}
\subcante{2. Estáticos: dominancia; EN; mixtas; minimax}{3,5}
\subcante{3. Dinámicos: secuenciales (ENPS) y repetidos}{4}
\apartadocante{II. Conceptos (II): incompleta y cooperativos}{7,5}
\subcante{1. Bayesianos (\autor{Harsanyi}): EBN}{2}
\subcante{2. EBP y señalización; criterio intuitivo}{2}
\subcante{3. Negociación; núcleo; valor de \autor{Shapley}}{3,5}
\apartadocante{III. Aplicaciones: especial referencia a subastas}{10}
\subcante{1. Aplicaciones en economía}{1,5}
\subcante{2. Diseño de mecanismos: revelación; límites}{2}
\subcante{3. Subastas: formatos, equivalencia, colusión}{6,5}
\apartadocante{Conclusión}{1}}
\emph{Ejemplo}: \textbf{I} (instalada) y \textbf{E} (entrante): guerra de precios $\to$ entrada $\to$ cada mes $\to$ costes ocultos $\to$ consorcio con O $\to$ pujan por el 5G.

% ==========================================================================
\bloque[2,5]{Introducción}
\begin{itemize}
\item \textbf{Enganche}: «Hoy por ti, mañana por mí»: ayudamos sin firmar nada a quien volveremos a ver; no somos mejores, cambia el \emph{juego}. Julio de 2021: subasta de 700~MHz (5G), 12 rondas, 1.010 M€ (salida 995,5 M€). \emph{¿Se puede predecir qué harán agentes interdependientes? ¿Y diseñar reglas para que revelen lo que saben?}
\item \textbf{Relevancia}: \emph{teórica} (lenguaje común de la micro; 6 Nobel: 1994, 1996, 2005, 2007, 2012, 2020); \emph{empírica} (oligopolio \vertema{3.A.19}, negociación salarial \vertema{3.A.27}, entrada \vertema{3.A.15}, guerras comerciales \vertema{3.B.7}, credibilidad del BC \vertema{3.A.36}); \emph{de política} (el Estado subasta deuda, espectro, renovables, CO$_2$ y persigue cárteles en licitaciones).
\item \textbf{Contextualización}: \autor{Cournot} (1838), \autor{Bertrand} (1883), \autor{Zermelo} (1913) $\to$ minimax de \autor{von Neumann} (1928) y \emph{Theory of Games} (1944) $\to$ 1950-75: \autor{Nash} (1950), \autor{Shapley} (1953), \autor{Selten} (1965), \autor{Harsanyi} (1967) $\to$ desde los 70, aplicación masiva; desde los 90, diseño de subastas y mercados. En competencia perfecta y monopolio no hace falta \vertema{3.A.16}\vertema{3.A.17}: sí con pocos agentes.
\item \textbf{Problemática}: ¿qué predecir cuando cada uno elige pensando en los demás? ¿Cuándo es creíble una amenaza? ¿Qué cambia con turnos, repetición, información privada? ¿Cómo repartir las ganancias de cooperar? ¿Qué formato de subasta conviene? ¿Por qué la cooperación del refrán es un problema en las licitaciones?
\item \textbf{Estructura}: \hyperref[rep:I]{I}, información completa (elementos, estáticos, dinámicos); \hyperref[rep:II]{II}, información incompleta, negociación y cooperativos; \hyperref[rep:III]{III}, aplicaciones y subastas. Eje: el \emph{concepto de equilibrio}, que se refina con cada ingrediente (tiempo, información privada, acuerdos); III le da la vuelta: elegir el juego que da el resultado deseado.
\end{itemize}
\begin{itemize}
\item \concepto{Teoría de juegos}: estudia situaciones de \concepto{interdependencia estratégica}: lo que obtiene cada agente depende de sus decisiones y de las de los demás, y cada uno lo sabe. \concepto{Juego}: su representación formal. La decisión con riesgo es el caso de 1 jugador frente a la naturaleza \vertema{3.A.10}.
\end{itemize}

% ==========================================================================
\bloque[9]{I. Principales conceptos (I): información completa\label{rep:I}}
\begin{ideaclave}
Con información completa, el EN es el concepto central: existe siempre (en mixtas) en juegos finitos, pero puede ser ineficiente, múltiple o apoyarse en amenazas no creíbles; el ENPS lo afina y la repetición sostiene la cooperación sin contratos.
\end{ideaclave}

\sub{I.1. Elementos, representación y clasificación\label{rep:I.1}}
\begin{itemize}
\item \emph{4 elementos} (\autor{Mas-Colell} et al., 1995): \concepto{jugadores} $i\in N$; \concepto{reglas} (quién mueve, cuándo, qué sabe, qué \concepto{acciones}); \concepto{resultados}; \concepto{pagos} $u_i$ (utilidad esperada vNM si hay azar \vertema{3.A.10}).
\item \concepto{Estrategia} $s_i\in S_i$: \emph{plan completo de acción}, qué hará en cada nodo en que pueda mover, incluso los que no se alcanzan $\neq$ acción. \concepto{Perfil} $s=(s_i,s_{-i})$. \concepto{Puras} o \concepto{mixtas} (al azar).
\item \concepto{Conocimiento común} (\autor{Aumann}, 1976): todos saben, todos saben que todos saben\ldots
\item \emph{Representación}: \concepto{forma extensiva} (árbol: nodos de decisión, ramas, nodos terminales con pagos; \concepto{conjunto de información}: nodos que el jugador no distingue, línea discontinua); \concepto{forma normal} (matriz). Cada extensiva tiene una normal, no al revés: la normal pierde el orden de las jugadas.
\item \emph{Clasificación}: \concepto{no cooperativos}/\concepto{cooperativos} (acuerdos vinculantes; «no cooperativo» = la cooperación, si la hay, se sostiene sola); \concepto{estáticos} (simultáneos)/\concepto{dinámicos} (\concepto{secuenciales}, una vez por turnos; \concepto{repetidos}, un \concepto{juego de etapa} finito o infinitas veces); información \concepto{perfecta}/\concepto{imperfecta} (sobre las \emph{jugadas}); \concepto{completa}/\concepto{incompleta} (sobre los \emph{pagos}: el \concepto{tipo}); suma cero/variable.
\end{itemize}
{\centering\small\renewcommand{\arraystretch}{1.1}\setlength{\tabcolsep}{2.5pt}%
\begin{tabular}{@{}>{\centering\arraybackslash}m{0.22\linewidth}>{\centering\arraybackslash}m{0.36\linewidth}>{\centering\arraybackslash}m{0.36\linewidth}@{}}
\toprule
\rowcolor{tceeazulpalido} & \textbf{Estáticos} & \textbf{Dinámicos}\\
\midrule
\textbf{Inf.\ completa} & \cellcolor{tceesalvia}\textcolor{tceeverde}{\textbf{EN}} (\hyperref[rep:I.2]{I.2}) & \cellcolor{tceesalvia}\textcolor{tceeverde}{\textbf{ENPS}} (\hyperref[rep:I.3]{I.3})\\
\textbf{Inf.\ incompleta} & \cellcolor{tceesalvia}\textcolor{tceedorado}{\textbf{EBN}} (\hyperref[rep:II.1]{II.1}) & \cellcolor{tceesalvia}\textcolor{tceedorado}{\textbf{EBP}} (\hyperref[rep:II.2]{II.2})\\
\rowcolor{tceesalmon}\textbf{Cooperativos} & \multicolumn{2}{c}{negociación de \autor{Nash}, núcleo, \autor{Shapley} (\hyperref[rep:II.3]{II.3})}\\
\bottomrule
\end{tabular}\par}
\begin{itemize}
\item \emph{Refinamientos}: cada concepto refina al anterior: todo ENPS es EN; todo EBP es EBN. Las subastas son juegos bayesianos.
\end{itemize}

\sub{I.2. Juegos estáticos: dominancia y equilibrio de \autor{Nash}\label{rep:I.2}}
\begin{itemize}
\item \concepto{Estrictamente dominante}: $u_i(s_i,s_{-i})>u_i(s_i',s_{-i})$ $\forall s_i'\neq s_i$, $\forall s_{-i}$. \concepto{Débilmente dominada}: $\geq$ siempre y $>$ en alguna.
\item \emph{Ejemplo: guerra de precios} (\textup{I} filas, E columnas):
\end{itemize}
{\centering\small\setlength{\tabcolsep}{4pt}%
\begin{tabular}{c|c|c|}
\multicolumn{1}{c}{} & \multicolumn{1}{c}{\textbf{Alto}} & \multicolumn{1}{c}{\textbf{Bajo}}\\
\cline{2-3}
\textbf{Alto} & $(5,5)$ & $(0,\underline{6})$\\
\cline{2-3}
\textbf{Bajo} & $(\underline{6},0)$ & \cellcolor{tceesalvia}$(\underline{1},\underline{1})$\\
\cline{2-3}
\end{tabular}\par}
\begin{itemize}
\item Bajar es dominante ($6>5$, $1>0$) $\Rightarrow$ (Bajo, Bajo), \emph{ineficiente en el sentido de \autor{Pareto}}: \concepto{dilema del prisionero} (\autor{Tucker}, 1950): cártel \vertema{3.A.19}, polizón \vertema{3.A.23}, guerras arancelarias \vertema{3.B.7}, clemencia \vertema{4.A.15}.
\item \emph{Eliminación iterativa} de estrictamente dominadas: el orden no importa (con débiles, sí); cada ronda exige un nivel más de conocimiento común. Pocas veces decide.
\item \concepto{Mejor respuesta} $R_i(s_{-i})=\arg\max_{s_i}u_i(s_i,s_{-i})$ (con acciones continuas, \concepto{función de reacción}). \concepto{Equilibrio de \autor{Nash}} (EN; 1950): cada estrategia es mejor respuesta a las de los demás:
\funcionclave{u_i(s_i^*,s_{-i}^*)\geq u_i(s_i,s_{-i}^*)\ \ \forall s_i,\forall i\ \Longleftrightarrow\ s_i^*\in R_i(s_{-i}^*)}
\item[] Desigualdad \emph{débil}. \emph{Expectativas que se confirman}. Todo EN sobrevive a la eliminación iterativa, no al revés. Cálculo: subrayar mejores respuestas; casilla con los 2 subrayados.
\item \concepto{Estrategia mixta} $\vect{\pi}_i\in\Delta(S_i)$: en el EN, cada jugador está \emph{indiferente} entre las puras que mezcla $\Rightarrow$ \emph{sus probabilidades dependen de los pagos del otro}.
\item \emph{Existencia} (\autor{Nash}, 1950): todo juego con nº finito de jugadores y estrategias puras tiene al menos un EN en mixtas (punto fijo de \autor{Kakutani}, 1941; \autor{Brouwer} en 1951). Continuas: EN en puras si estrategias compactas y convexas y pagos continuos y cuasicóncavos (\autor{Cournot}).
\item \emph{Ejemplo: la campaña} (\textup{I} quiere coincidir, E evitarlo): sin EN en puras; mixtas al 50\,\%. \emph{Clásicos}: pares o nones; \emph{inspección} (si vigilar se encarece, no se vigila menos: el trabajador se escaquea más).
\end{itemize}
\graficorepaso[0.42\linewidth]{inspeccion}{Inspección: mejores respuestas y EN en mixtas ($\pi_V^*=1/2$, $\pi_T^*=3/4$)}
\begin{itemize}
\item \emph{Implicaciones}: predicción mínima (si todos la conocen, ha de ser EN); base de los refinamientos; acciones continuas: corte de las funciones de reacción $=$ \autor{Cournot} (1838); \autor{Bertrand}: $p=\mathrm{CMg}$ \vertema{3.A.19}.
\item \concepto{Teorema minimax} (\autor{von Neumann}, 1928): 2 jugadores, suma cero:
\begin{formula}
\centering $\max_{\vect{\pi}_1}\min_{\vect{\pi}_2}U_1(\vect{\pi}_1,\vect{\pi}_2)=\min_{\vect{\pi}_2}\max_{\vect{\pi}_1}U_1(\vect{\pi}_1,\vect{\pi}_2)$
\end{formula}
\item[] \concepto{Valor del juego}; minimax = EN. No confundir con maximin de \autor{Wald} (sin interdependencia) \vertema{3.A.10}. \emph{Evidencia}: penaltis (\autor{Palacios-Huerta}, 2003: $>$1.400): igualan probabilidad de marcar por cada lado.
\item \emph{Valoración: 4 límites del EN}: i) puede no existir en puras; ii) \emph{varios}: \concepto{juegos de coordinación}, \concepto{puntos focales} (\autor{Schelling}, 1960: Puerta del Sol); ejemplo: estándar técnico = \concepto{batalla de los sexos} (\autor{Luce}-\autor{Raiffa}, 1957); amenazas de guerra = \concepto{juego de la gallina} (España-BCE, 2012); iii) no \emph{eficiente en el sentido de \autor{Pareto}}; iv) \emph{amenazas no creíbles} en dinámicos.
\end{itemize}

\sub{I.3. Juegos dinámicos: secuenciales y repetidos\label{rep:I.3}}
\begin{itemize}
\item \emph{Secuenciales: problema nuevo}, la \concepto{credibilidad}. \emph{Ejemplo: entrada de E}: E entra o no; si entra, I acomoda o lucha. 2 EN: (Entrar, Acomodar) y (No entrar, Luchar si entra), este con amenaza \emph{no creíble} ($1>-1$): «Perro ladrador, poco mordedor».
\end{itemize}
\graficorepaso[0.85\linewidth]{juego-entrada-compromiso}{Entrada de E y compromiso de \textup{I}: invertir antes en capacidad hace creíble luchar (en rojo, el ENPS)}
\begin{itemize}
\item \concepto{Inducción hacia atrás}: un \emph{procedimiento}: última decisión $\to$ sustituir el subárbol $\to$ repetir hacia la raíz. Juego finito de información perfecta: equilibrio en puras, único si no hay empates (\autor{Zermelo} 1913; \autor{Kuhn} 1953).
\item \concepto{Subjuego}: parte del árbol desde un nodo con conjunto de información unitario, con todos sus sucesores, sin cortar conjuntos de información.
\item \concepto{Equilibrio de \autor{Nash} perfecto en subjuegos} (ENPS; \autor{Selten}, 1965): induce un EN en \emph{todo} subjuego, también fuera de la senda: \concepto{racionalidad secuencial}. Refinamiento del EN. Info.\ perfecta: ENPS = inducción hacia atrás.
\item \emph{Compromiso} (\autor{Schelling}, 1960): reducir las propias opciones da más. Ej.: \textup{I} invierte en capacidad (hundido 2,5): si E entra, acomodar $-1{,}5$ < luchar $-1$ $\Rightarrow$ E no entra, I gana $4-2{,}5=1{,}5>1$ (\autor{Spence} 1977, \autor{Dixit} 1980) \vertema{3.A.15}. «Quemar las naves» (\autor{Cortés}, 1519).
\item \emph{Inconsistencia temporal}: conjetura de \autor{Coase} (1972) \vertema{3.A.17}; \autor{Kydland}-\autor{Prescott} (1977), \autor{Barro}-\autor{Gordon} (1983) \vertema{3.A.36}.
\item \emph{Extensiones}: juegos en 2 etapas (capacidad/localización $\to$ precios o cantidades) \vertema{3.A.18}\vertema{3.A.19}; \autor{Stackelberg} = ENPS de \autor{Cournot} por turnos; \autor{Rubinstein}.
\item \emph{Evidencia}: ciempiés (\autor{Rosenthal}, 1981; \autor{McKelvey}-\autor{Palfrey}, 1992): casi nadie para al principio. \emph{Valoración}: sin información perfecta apenas hay subjuegos $\Rightarrow$ creencias (EBP).
\item \emph{Repetidos}. Ej.: \textup{I} y E se encuentran cada mes: «Hoy por ti, mañana por mí». Pagos descontados con $\beta\equiv1/(1+\rho)$. \concepto{Principio de la desviación única}: ENPS sii nadie gana desviándose en \emph{una sola} etapa tras cualquier historia.
\item \emph{Horizonte finito} con EN de etapa único: único ENPS = repetirlo (inducción hacia atrás). \concepto{Paradoja de la cadena de tiendas} (\autor{Selten}, 1978; 20 ciudades): acomoda en todas \vertema{3.A.15}.
\item \emph{Horizonte infinito} (o final incierto): \concepto{incondicionada} (EN de etapa siempre: ENPS, no saca del dilema) frente a \concepto{condicionadas}: \concepto{gatillo} (\autor{Friedman}, 1971): cooperar mientras el otro coopere; si se desvía, EN para siempre. $u^{\mathrm C}/(1-\beta)\geq u^{\mathrm D}+\beta u^{\mathrm N}/(1-\beta)$:
\funcionclave{\beta\geq\frac{u^{\mathrm D}-u^{\mathrm C}}{u^{\mathrm D}-u^{\mathrm N}}}
\item[] Numerador: ganancia hoy; denominador: pérdida cada etapa después. Castigo creíble $\Rightarrow$ ENPS. Ej.: $(6-5)/(6-1)=1/5$; cada mes, $\beta\approx1$.
\item \emph{Otras}: \emph{tit for tat} (ENPS solo con $\beta=1/5$ exacto); \emph{tit for two tats} (no: ni EN); \emph{win-stay, lose-shift} (sí si $\beta\geq1/4$ y $2u^{\mathrm C}>u^{\mathrm D}+u^{\mathrm N}$).
\item \concepto{Teorema de tradición oral} (\emph{folk theorem}): con infinitas repeticiones, \emph{cualquier} pago medio factible por encima del castigo es ENPS si $\beta\to1$ (\autor{Friedman} 1971, reversión al EN; \autor{Fudenberg}-\autor{Maskin} 1986, minimax). Ej.: entre 1 y 5.
\item \emph{Implicaciones}: colusión tácita \vertema{3.A.19}, reputación \vertema{3.A.36}, acuerdos comerciales \vertema{3.B.7}; también cárteles en licitaciones. Más cooperación con más paciencia, más frecuencia, detección rápida. \emph{Reputación} con tipos «irracionales» (\autor{Kreps}-\autor{Milgrom}-\autor{Roberts}-\autor{Wilson}, 1982) deshace la paradoja.
\item \emph{Evidencia}: torneos de \autor{Axelrod} (1984): gana \emph{tit for tat}; con errores, \emph{win-stay} (\autor{Nowak}-\autor{Sigmund}, 1993); \autor{Dal Bó} (2005). \emph{Valoración}: permite demasiado $\Rightarrow$ \emph{selección de equilibrios}.
\end{itemize}

% ==========================================================================
\bloque[7,5]{II. Principales conceptos (II): información incompleta, negociación y cooperativos\label{rep:II}}
\begin{ideaclave}
La información privada se trata con tipos, probabilidades y creencias (EBN, EBP); con acuerdos vinculantes, la pregunta pasa a ser cómo repartir lo que el grupo puede conseguir (\autor{Nash}, núcleo, \autor{Shapley}).
\end{ideaclave}

\sub{II.1. Juegos bayesianos: equilibrio bayesiano de \autor{Nash}\label{rep:II.1}}
\begin{itemize}
\item \emph{Truco de \autor{Harsanyi}} (1967): la \concepto{Naturaleza} elige el \concepto{tipo} $t_i$ de cada uno según una \concepto{distribución a priori} común; cada uno ve el suyo $\Rightarrow$ información incompleta $\to$ \emph{imperfecta}.
\item \emph{Estrategia}: \emph{regla} $s_i(t_i)$ según el tipo (los rivales tienen que anticipar qué haría cada tipo).
\item \concepto{Equilibrio bayesiano de \autor{Nash}} (EBN): cada tipo maximiza su pago \emph{esperado}, dadas creencias y estrategias: EN del juego en que cada tipo es un jugador:
\begin{formula}
\centering $\mathbb{E}_{t_{-i}}\big[u_i(s_i^*(t_i),s_{-i}^*(t_{-i});t_i)\,|\,t_i\big]\geq\mathbb{E}_{t_{-i}}\big[u_i(s_i,s_{-i}^*(t_{-i});t_i)\,|\,t_i\big]$
\end{formula}
\item \emph{Ejemplo}: E entra mientras \textup{I} decide si amplía; I de coste alto ($\pi$) nunca amplía, de coste bajo siempre (dominantes). E: $\pi\cdot1+(1-\pi)(-1)=2\pi-1\geq0$ $\Rightarrow$ entra si $\pi\geq1/2$.
\end{itemize}
\graficorepaso[0.8\linewidth]{arbol-harsanyi}{Entrada sin conocer el coste de \textup{I}: la Naturaleza elige el tipo; línea discontinua, nodos que E no distingue}
\begin{itemize}
\item \emph{Implicaciones}: se responde a lo que se \emph{espera}; subastas (tipo $=v_i$, estrategia $b(v_i)$) y selección adversa \vertema{3.A.13}. \emph{Extensiones}: \autor{Cournot} con costes privados; \emph{purificación} (\autor{Harsanyi}, 1973). \emph{Valoración}: la \concepto{doctrina de \autor{Harsanyi}} (a priori común) es fuerte.
\end{itemize}

\sub{II.2. Equilibrio bayesiano perfecto y señalización\label{rep:II.2}}
\begin{itemize}
\item \emph{Idea}: si el informado mueve primero, su acción revela su tipo. Ej.: el precio de \textup{I} es una señal; E actualiza con la \concepto{regla de \autor{Bayes}}: «Cuando el río suena, agua lleva» (precio límite, \autor{Milgrom}-\autor{Roberts} 1982 \vertema{3.A.15}):
\begin{formula}
\centering $\mu(\text{precio})=\dfrac{\pi\Pr(\text{precio}|\text{alto})}{\pi\Pr(\text{precio}|\text{alto})+(1-\pi)\Pr(\text{precio}|\text{bajo})}$
\end{formula}
\item \emph{Modelo} (\autor{Spence}, 1973, como juego \vertema{3.A.13}): tipos A ($\nu$) y B; $w_{\mathrm A}>w_{\mathrm B}$; educación $y$ improductiva, $c_{\mathrm A}(y)<c_{\mathrm B}(y)$, $c_{\mathrm A}'<c_{\mathrm B}'$ (cruce único); $w(y)=\mu(y)w_{\mathrm A}+[1-\mu(y)]w_{\mathrm B}$.
\item \concepto{Equilibrio bayesiano perfecto} (EBP): estrategias + \concepto{creencias} con 3 condiciones: i) \emph{racionalidad secuencial} dadas las creencias; ii) \emph{\autor{Bayes} en la senda}; iii) \emph{fuera de la senda}, libres (EBP débil) o con requisitos.
\item \concepto{Separadores}: A elige $\hat y$, B 0; $\mu(\hat y)=1$, $\mu(0)=0$, si
\begin{formula}
\centering $c_{\mathrm A}(\hat y)\leq w_{\mathrm A}-w_{\mathrm B}\leq c_{\mathrm B}(\hat y)$
\end{formula}
\item \concepto{Agrupadores}: misma $y$, $\mu=\nu$, $\overline w=\nu w_{\mathrm A}+(1-\nu)w_{\mathrm B}$; con creencias pesimistas fuera de la senda.
\item \concepto{Criterio intuitivo} (\autor{Cho}-\autor{Kreps}, 1987): quien elige lo que B no elegiría nunca y A sí $\Rightarrow$ es A. Elimina agrupadores y deja el separador más barato: $c_{\mathrm B}(\hat y^*)=w_{\mathrm A}-w_{\mathrm B}$. \emph{Clásico}: \emph{cerveza o quiche}.
\item \emph{Implicaciones}: las creencias fuera del equilibrio sostienen el resultado; la señal cuesta eficiencia. \emph{Refinamientos}: \concepto{equilibrio secuencial} (\autor{Kreps}-\autor{Wilson}, 1982); \concepto{mano temblorosa} (\autor{Selten}, 1975). \emph{Valoración}: completa la tabla, pero vuelve la multiplicidad.
\end{itemize}

\sub{II.3. Negociación y juegos cooperativos\label{rep:II.3}}
\begin{itemize}
\item \emph{Idea}: indeterminación de \autor{Edgeworth} \vertema{3.A.21} y \autor{Bowley} \vertema{3.A.17}; \concepto{programa de \autor{Nash}}: los caminos cooperativo y no cooperativo llegan al mismo sitio. Ej.: acceso de E a la red de \textup{I}.
\item \emph{Supuestos}: $\mathcal{U}$ convexo y compacto; \concepto{punto de desacuerdo} $(u_1^0,u_2^0)$.
\item \emph{Axiomático} (\autor{Nash}, 1950): eficiencia, simetría, invariancia afín, independencia de alternativas irrelevantes $\Rightarrow$ máx.\ producto de ganancias ($\gamma=1/2$); \concepto{generalizada}, \concepto{poder de negociación} $\gamma$:
\funcionclave{\optim{\max}{u_1,u_2}{(u_1-u_1^0)^{\gamma}(u_2-u_2^0)^{1-\gamma}}{(u_1,u_2)\in\mathcal{U}}}
\item \emph{Estratégico} (\autor{Rubinstein}, 1982): ofertas alternadas por 1 €, descuentos $\beta_1,\beta_2$; único ENPS, acuerdo inmediato:
\begin{formula}
\centering $x^*=\dfrac{1-\beta_2}{1-\beta_1\beta_2}$;\quad $\Delta\to0$: $\gamma=\dfrac{\rho_2}{\rho_1+\rho_2}$
\end{formula}
\item[] Gana el más paciente (\autor{Binmore}-\autor{Rubinstein}-\autor{Wolinsky}, 1986): poder = \emph{paciencia relativa}. Aplicaciones: monopolio bilateral \vertema{3.A.17}, salarios \vertema{3.A.27}, cártel \vertema{3.A.19}. \emph{Evidencia}: \concepto{ultimátum} (\autor{Güth} et al., 1982): ofertas cercanas a la mitad, rechazo de las bajas \vertema{3.A.1}.
\item \emph{Cooperativos}: acuerdos vinculantes y \concepto{utilidad transferible}; \concepto{coaliciones} $T\subseteq N$; \concepto{función característica} $v(T)$; \concepto{imputación}: eficiente ($\sum x_i=v(N)$) e individualmente racional ($x_i\geq v(\{i\})$).
\item \concepto{Núcleo} (\autor{Gillies}, 1959): imputaciones que ninguna coalición \concepto{bloquea}:
\begin{formula}
\centering $\mathcal{C}(v)=\{x:\sum_{N}x_i=v(N),\ \sum_{T}x_i\geq v(T)\ \forall T\}$
\end{formula}
\item \emph{Ejemplo}: \textup{I} vende su red (20 solo); E la valora en 60, O en 100. Núcleo: $x_2=0$, $x_1\in[60,100]$, $x_3=100-x_1$: la red va a O a un precio entre la 2.ª y la 1.ª valoración (como una subasta inglesa).
\end{itemize}
\graficorepaso[0.5\linewidth]{nucleo}{Imputaciones y núcleo del juego de \textup{I} (vendedor), E y O}
\begin{itemize}
\item \emph{Límites}: puede ser \emph{vacío} (mayoría por 100); \autor{Bondareva}-\autor{Shapley}: no vacío sii juego \emph{equilibrado}; rara vez un punto; se encoge al walrasiano (\autor{Debreu}-\autor{Scarf}, 1963) \vertema{3.A.21}.
\item \concepto{Valor de \autor{Shapley}} (1953): 4 axiomas (eficiencia, simetría, jugador nulo, aditividad) $\Rightarrow$ regla única:
\begin{formula}
\centering $\phi_i(v)=\displaystyle\sum_{T\subseteq N\setminus\{i\}}\frac{|T|!\,(n-|T|-1)!}{n!}\big[v(T\cup\{i\})-v(T)\big]$
\end{formula}
\item[] \emph{Aportación marginal esperada} con orden de llegada aleatorio. Ej.: $\phi=(66\tfrac23,\,6\tfrac23,\,26\tfrac23)$: \emph{no} está en el núcleo (da algo a E). Siempre existe y es único; en el núcleo si el juego es \emph{convexo} (\autor{Shapley}, 1971). Usos: poder de voto (\autor{Shapley}-\autor{Shubik}, 1954) \vertema{3.A.24}, costes comunes, cárteles, aprendizaje automático.
\end{itemize}

% ==========================================================================
\bloque[10]{III. Aplicaciones, con especial referencia a las subastas\label{rep:III}}
\begin{ideaclave}
El diseñador elige las reglas: basta buscar mecanismos en que decir la verdad sea óptimo; la información privada tiene un coste que el diseño reduce pero no elimina; en subastas, con valores privados independientes, el formato importa menos de lo que parece: lo decisivo es atraer pujadores y evitar la colusión.
\end{ideaclave}

\sub{III.1. Aplicaciones en economía\label{rep:III.1}}
\begin{itemize}
\item \emph{OI}: \autor{Cournot}, \autor{Bertrand}, \autor{Stackelberg}, colusión tácita \vertema{3.A.19}; diferenciación por etapas \vertema{3.A.18}; entrada, capacidad, precio límite \vertema{3.A.15}.
\item \emph{Información asimétrica}: señalización, autoselección, principal-agente \vertema{3.A.13}; regulación \vertema{3.A.17}. \emph{Trabajo}: negociación \vertema{3.A.27}.
\item \emph{Sector público}: polizón, \autor{Clarke}-\autor{Groves} \vertema{3.A.23}; voto estratégico, \autor{Gibbard}-\autor{Satterthwaite} \vertema{3.A.24}; imposición óptima \vertema{4.B.11}.
\item \emph{Comercio}: guerra arancelaria (\autor{Johnson}, 1953) \vertema{3.B.7}; \autor{Brander}-\autor{Spencer} (1985) \vertema{3.B.8}. \emph{Macro}: inconsistencia temporal \vertema{3.A.36}; ataques especulativos \vertema{3.B.17}.
\item \emph{Competencia}: clemencia \vertema{4.A.15}. \emph{Diseño de mercados}: aceptación diferida (\autor{Gale}-\autor{Shapley}, 1962; \autor{Roth}). \emph{Evolutiva} (\autor{Maynard Smith}-\autor{Price}, 1973) y conductual \vertema{3.A.1}.
\end{itemize}

\sub{III.2. Diseño de mecanismos y principio de revelación\label{rep:III.2}}
\begin{itemize}
\item \emph{Idea}: elegir las reglas (\autor{Hurwicz}, 1960, 1972; debate del cálculo socialista \vertema{3.A.21}). Ej.: el Gobierno no sabe cuánto valen las licencias.
\item \concepto{Mecanismo}: mensajes $\to$ asignación y pagos (tipos $\to$ mensajes $\to$ resultado); induce un juego bayesiano. \concepto{Directo}: el mensaje es el tipo.
\item \concepto{Principio de revelación} (\autor{Gibbard} 1973; \autor{Dasgupta}-\autor{Hammond}-\autor{Maskin} 1979; \autor{Myerson} 1979): todo EBN equivale a un mecanismo directo compatible con los incentivos en que todos dicen la verdad. Basta buscar con $\mathrm{PC}_i$ (participación) e $\mathrm{IC}_i$ (compatibilidad de incentivos).
\item \emph{Clásico}: juicio de \autor{Salomón}: funcionó por sorpresa, no en equilibrio.
\item \concepto{VCG} (\autor{Vickrey} 1961, \autor{Clarke} 1971, \autor{Groves} 1973): pagar el daño que se causa a los demás $\Rightarrow$ verdad \emph{dominante} y decisión eficiente; caso más sencillo, 2.º precio.
\item \emph{Implicaciones}: hay que dejar \concepto{renta informativa} \vertema{3.A.13}\vertema{3.A.17}\vertema{4.B.11}. \emph{Implementación} (\autor{Maskin}, 1999): en todos los equilibrios, monotonía.
\item \emph{Límites}: \autor{Gibbard}-\autor{Satterthwaite} ($\geq3$ alternativas: no manipulable $\Rightarrow$ dictatorial; VCG necesita dinero) \vertema{3.A.24}; \autor{Myerson}-\autor{Satterthwaite} (1983): intercambio bilateral, ningún mecanismo eficiente, IC, voluntario y sin déficit; el diseñador ha de poder comprometerse.
\end{itemize}

\sub{III.3. Subastas\label{rep:III.3}}
\begin{itemize}
\item \concepto{Subasta}: mecanismo de asignación y precio a partir de las pujas; la competencia revela valoraciones y reduce rentas de información (\autor{Klemperer}, 2004).
\item \emph{Formatos}: \concepto{abiertos}: \concepto{inglesa} (ascendente) y \concepto{holandesa} (descendente; Aalsmeer, lonjas); \concepto{sobre cerrado}: \concepto{primer precio} (paga su puja) y \concepto{segundo precio} o de \autor{Vickrey} (paga la 2.ª). Abierta $\neq$ secuencial.
\item \emph{Valores}: \concepto{privados} ($v_i$ propio; VPI si independientes); \concepto{común} (igual para todos, nadie lo sabe: petróleo); \concepto{afiliados} (\autor{Milgrom}-\autor{Weber}, 1982). Juego bayesiano: estrategia $b(v_i)$, solución EBN.
\end{itemize}
{\centering\small\renewcommand{\arraystretch}{1.1}\setlength{\tabcolsep}{2.5pt}%
\begin{tabular}{@{}>{\raggedright\arraybackslash}m{0.24\linewidth}>{\raggedright\arraybackslash}m{0.34\linewidth}>{\raggedright\arraybackslash}m{0.36\linewidth}@{}}
\toprule
\rowcolor{tceeazulpalido}\textbf{Formato} & \textbf{Gana y paga} & \textbf{Puja (VPI)}\\
\midrule
\textcolor{tceeverde}{\textbf{Inglesa}} & último; precio al que se retira el penúltimo & \textcolor{tceeverde}{hasta $v_i$ (= 2.º precio)}\\
\textcolor{tceeverde}{\textbf{2.º precio}} & puja más alta; paga la 2.ª & \textcolor{tceeverde}{$b=v_i$ (dominante)}\\
\textcolor{tceedorado}{\textbf{1.\textsuperscript{er} precio}} & puja más alta; paga la suya & \textcolor{tceedorado}{sombrear: $\frac{n-1}{n}v_i$}\\
\textcolor{tceedorado}{\textbf{Holandesa}} & primero que acepta & \textcolor{tceedorado}{= 1.\textsuperscript{er} precio (siempre)}\\
\midrule
\multicolumn{3}{c}{\cellcolor{tceesalvia}\emph{Equivalencia}: los 4 eficientes; ingreso $\frac{n-1}{n+1}$ con $U[0,1]$}\\
\bottomrule
\end{tabular}\par}
\begin{itemize}
\item \emph{Modelo de referencia (VPI)}: 1 objeto, $n$ pujadores \emph{neutrales al riesgo} y \emph{simétricos}, $v_i$ independientes de $F$ continua ($U[0,1]$); $v_0=0$, sin reserva.
\item \emph{2.º precio}: $b(v_i)=v_i$ \emph{débilmente dominante} (\autor{Vickrey}, 1961): la puja decide \emph{si} gana, no \emph{cuánto} paga; eficiente. Ej.: $v_{\mathrm I}=10$, $v_{\mathrm E}=8$, $v_{\mathrm O}=6$: I gana y paga 8. \emph{Inglesa}: equivalente con valores privados.
\item \emph{1.\textsuperscript{er} precio}: dilema (más puja $\uparrow$ prob.\ de ganar, $\downarrow$ ganancia): $\max_{b_i}(v_i-b_i)\Pr(\text{ganar})$; con $U[0,1]$:
\funcionclave{b(v_i)=\frac{n-1}{n}\,v_i=\mathbb{E}\big[Y_1\,|\,Y_1<v_i\big]}
\item[] Puja lo que espera pagar en 2.º precio; $\uparrow n$ $\Rightarrow$ sombrea menos. Ej. ($U[0,12]$, $n=3$): 6,67, 5,33 y 4. \emph{Holandesa}: estratégicamente equivalente al 1.\textsuperscript{er} precio con cualquier valor.
\end{itemize}
\pareja{pujas}{Pujas de equilibrio con VPI uniformes: $b=v$ y $b=\frac{n-1}{n}v$}{multiunidad}{Varias unidades: precio uniforme frente a discriminatorio}
\begin{itemize}
\item \concepto{Teorema de equivalencia de ingresos} (\autor{Vickrey} 1961; \autor{Myerson} 1981; \autor{Riley}-\autor{Samuelson} 1981): neutrales, simétricos, VPI; todo formato en que i) gana quien más valora y ii) el tipo más bajo espera 0 da el \emph{mismo ingreso esperado}:
\begin{formula}
\centering $\mathbb{E}[I^{\mathrm{1P}}]=\tfrac{n-1}{n}\mathbb{E}[v_{(1)}]=\mathbb{E}[v_{(2)}]=\mathbb{E}[I^{\mathrm{2P}}]=\tfrac{n-1}{n+1}$
\end{formula}
\item[] Intuición (envolvente): lo que espera ganar cada tipo depende solo de con qué probabilidad gana. Ej.: 8 frente a 6,67; en promedio, $12\cdot\frac24=6$. También «todos pagan».
\item \emph{Subasta óptima} (\autor{Myerson}; \autor{Riley}-\autor{Samuelson}, 1981): formato estándar + \concepto{precio de reserva} con $\psi(p^R)=v_0$; \concepto{valor virtual} $\psi(v)=v-\frac{1-F(v)}{f(v)}$ (IMg del monopolista \vertema{3.A.17}). $U[0,1]$: $\psi=2v-1$, $p^R=1/2$ $\forall n$; con 2: $1/3\to5/12$. Sacrifica eficiencia. \autor{Bulow}-\autor{Klemperer} (1996): un pujador más ($1/2$) vale más que la reserva.
\item \emph{Aversión al riesgo}: 1.\textsuperscript{er} precio recauda más (\autor{Holt} 1980; \autor{Maskin}-\autor{Riley} 1984). \emph{Asimetría}: sin ordenación general (\autor{Maskin}-\autor{Riley}, 2000).
\item \concepto{Maldición del ganador}: con valor común gana quien más sobrestimó; racional: pujar \emph{por debajo} de la estimación, más con más rivales (\autor{Wilson}, 1969, 1977). «No es oro todo lo que reluce». \emph{Clásicos}: golfo de México (\autor{Capen}-\autor{Clapp}-\autor{Campbell}, 1971); jarra de monedas (\autor{Bazerman}-\autor{Samuelson}, 1983).
\item \concepto{Principio de vinculación} (\autor{Milgrom}-\autor{Weber}, 1982): valores afiliados $\Rightarrow$ ingreso esperado inglesa $\geq$ 2.º precio $\geq$ 1.\textsuperscript{er} precio = holandesa; al vendedor le conviene revelar información.
\item \emph{Evidencia}: \autor{Kagel}-\autor{Levin} (1986, 1993): maldición en el laboratorio; en 1.\textsuperscript{er} precio se puja de más. Reserva en acción: renovables y 26~GHz en 2022 (datos clave).
\item \emph{Valoración} (\autor{Klemperer}, 2002): importa \emph{atraer pujadores} y \emph{evitar la colusión}. 3G (2000): RU, 5 licencias para 4 instalados. Ascendentes desaniman a los débiles; híbrida \emph{anglo-holandesa}.
\item \emph{Varias unidades}: \concepto{precio uniforme} (todos pagan la última aceptada $\Rightarrow$ \concepto{reducción de la demanda}) frente a \concepto{discriminatorio} (cada uno sus pujas); ninguno domina. Objetos distintos con complementariedades: \concepto{SMRA} (\autor{Milgrom}, \autor{Wilson}, \autor{McAfee}; regla de actividad; FCC 1994: 10 licencias, 47 rondas, 617 M\$); \concepto{combinatorias}, reloj, incentivos (FCC 2016-17).
\item \emph{España y UE}: Tesoro, «convencional corregida» (sobre el precio medio pagan el medio; entre medio y mínimo, lo ofrecido) \vertema{4.B.26}; espectro 3,6-3,8~GHz (2018), 700~MHz (2021) \vertema{4.A.12}; renovables (REER: paga lo ofrecido, con reserva) \vertema{4.A.5}; derechos de emisión UE: ronda única, sobre cerrado, precio uniforme \vertema{4.A.6}.
\item \concepto{Colusión} (\emph{bid rigging}): \concepto{ofertas de cobertura} y \concepto{rotación de ofertas} (\autor{Asker}, 2010): «Hoy gano yo, mañana tú». Con 2, el que espera: $G\beta/(1-\beta^2)\geq G$:
\begin{formula}
\centering $\dfrac{\beta}{1-\beta^2}\geq1\ \Longleftrightarrow\ \beta\geq\dfrac{\sqrt5-1}{2}\approx0{,}62$;\quad con 3: $\dfrac{\beta^2}{1-\beta^3}\geq1$, $\beta\gtrsim0{,}755$
\end{formula}
\end{itemize}
\graficorepaso[0.68\linewidth]{colusion-tiempo}{Gatillo y colusión: lo que se gana desviándose (rojo) frente a lo que se pierde con el castigo (azul)}
\begin{itemize}
\item[] Subastas cada 5 años al 10\,\%: $\beta=1/1{,}1^5\approx0{,}62$.
\item \emph{Remedios} (que no haya «mañana»): sobre cerrado; no publicar pujas; más pujadores; precio de reserva; clemencia \vertema{4.A.15}. \emph{Clásico}: Alemania 1999 (Mannesmann 18,18 y 20; $18{,}18\times1{,}1\approx20$; 2 rondas).
\end{itemize}

% ==========================================================================
\bloque[1]{Conclusión}
\begin{itemize}
\item \textbf{Recapitulación}: EN (siempre en mixtas; ineficiente o múltiple) $\to$ ENPS (credibilidad, compromiso; repetición sostiene la cooperación) $\to$ EBN y EBP (creencias, señales) $\to$ \autor{Nash}, núcleo, \autor{Shapley} $\to$ diseño de mecanismos: equivalencia; valor común, entrada y colusión deciden. I y E han recorrido todo el camino.
\item \textbf{Relevancia}: Nobel 2020, \autor{Milgrom} y \autor{Wilson}: maldición del ganador, vinculación, formatos para muchos objetos (eficiencia, no máxima recaudación); espectro en medio mundo, también en España.
\item \textbf{Extensiones}: oligopolio \vertema{3.A.19}, entrada \vertema{3.A.15}, información asimétrica \vertema{3.A.13}, elección colectiva \vertema{3.A.24}, credibilidad \vertema{3.A.36}, competencia \vertema{4.A.15}, conductual \vertema{3.A.1}.
\item \textbf{Opinión}: demasiados equilibrios y racionalidad exigente (conocimiento común); pero su mayor éxito es práctico: se juzga por cómo funcionan las instituciones que diseña (teoría + experimentos + experiencia).
\item \textbf{Idea final}: «Hoy por ti, mañana por mí»: la repetición sostiene la cooperación, para bien (reputación) y para mal («hoy gano yo, mañana tú»). Una buena subasta se diseña para que no haya «mañana» que repartir.
\end{itemize}

% ==========================================================================
\bloque{Autores y fechas}
\hito{1838}{\autor{Cournot}: duopolio (caso del EN)}
\hito{1928}{\autor{von Neumann}: teorema minimax}
\hito{1944}{\autor{von Neumann}-\autor{Morgenstern}: \emph{Theory of Games}}
\hito{1950}{\autor{Nash}: equilibrio y negociación; \autor{Tucker}: prisionero}
\hito{1953}{\autor{Shapley}: valor; \autor{Kuhn}: inducción hacia atrás}
\hito{1960}{\autor{Schelling}: puntos focales, compromiso; \autor{Hurwicz}}
\hito{1961}{\autor{Vickrey}: subasta de 2.º precio, equivalencia}
\hito{1965}{\autor{Selten}: perfección en subjuegos (ENPS)}
\hito{1967}{\autor{Harsanyi}: información incompleta}
\hito{1971}{\autor{Friedman}: gatillo; \autor{Capen} et al.: maldición}
\hito{1973}{\autor{Spence}: señalización; \autor{Gibbard}}
\hito{1978}{\autor{Selten}: paradoja de la cadena de tiendas}
\hito{1979}{\autor{Myerson}: principio de revelación}
\hito{1981}{\autor{Myerson}; \autor{Riley}-\autor{Samuelson}: subasta óptima}
\hito{1982}{\autor{Rubinstein}; \autor{Milgrom}-\autor{Weber}; KMRW}
\hito{1984}{\autor{Axelrod}: torneos (\emph{tit for tat})}
\hito{1986}{\autor{Fudenberg}-\autor{Maskin}: \emph{folk theorem}}
\hito{1987}{\autor{Cho}-\autor{Kreps}: criterio intuitivo}
\hito{1996}{\autor{Bulow}-\autor{Klemperer}: un pujador más}
\hito{2002}{\autor{Klemperer}: lo que de verdad importa}

\bloque{Datos clave}
\begin{itemize}
\item Subasta de 700~MHz (jul.\ 2021): \dato{1.010 M€ tras 12 rondas (salida 995,5 M€)}{Ministerio, 2021}.
\item 3,6-3,8~GHz (2018): \dato{437,6 M€, 34 rondas, $>$4 veces la salida}{Ministerio, 2018}.
\item Tesoro, 2025: \dato{46 subastas, 80\,\% de la financiación a medio y largo plazo}{Tesoro, 2026}.
\item 1.ª subasta de renovables (ene.\ 2021): \dato{9.700~MW ofertados para 3.034; 24,47 €/MWh FV, 25,31 eólica}{MITECO, 2021}.
\item 4.ª subasta (nov.\ 2022): \dato{45,5 de 3.300~MW}{MITECO, 2022}; 26~GHz (dic.\ 2022): \dato{9 de 12 estatales y 1 de 38 autonómicas, 36,2 M€}{Ministerio, 2022}.
\item 3G (2000): \dato{RU y Alemania $>$600 €/hab.; Países Bajos 170; Suiza 20}{\autor{Klemperer}, 2002}.
\item FCC, SMRA: \dato{1994: 617 M\$; en 20 años, $>$120.000 M\$}{Nobel, 2020}.
\item CNMC (jul.\ 2022): \dato{203,6 M€ a 6 constructoras, 25 años de licitaciones}{CNMC, 2022}.
\end{itemize}

% ==========================================================================
\bloque{Preguntas del tribunal}
\begin{enumerate}
\item \textbf{¿Diferencia entre información completa y perfecta?} Completa: pagos conocidos; perfecta: jugadas anteriores conocidas. El prisionero: completa e imperfecta. \autor{Harsanyi} convierte la incompleta en imperfecta.
\item \textbf{¿Por qué el ENPS y no el EN en la entrada?} El EN (No entrar, Luchar) descansa en una amenaza no creíble; el ENPS exige EN en todo subjuego.
\item \textbf{¿De qué dependen las probabilidades de un EN en mixtas?} De los pagos del \emph{otro} jugador (indiferencia).
\item \textbf{¿Se puede cooperar en un dilema repetido?} Infinito y $\beta$ alto: sí (gatillo, \emph{folk theorem}); finito con EN único: no.
\item \textbf{¿Está el valor de \autor{Shapley} en el núcleo?} No siempre; sí en juegos convexos. En el ejemplo, no.
\item \textbf{¿Qué formato de subasta recauda más?} Con VPI, neutrales y simétricos, todos igual; con afiliados, la inglesa; con aversión al riesgo, 1.\textsuperscript{er} precio. En la práctica, atraer pujadores y evitar la colusión.
\end{enumerate}

\begin{trampa}
\begin{itemize}
\item Desigualdad del EN \emph{débil}; mixtas: dependen de los pagos del otro (encarecer la vigilancia no cambia cuánto se vigila).
\item Eliminación iterativa: con estrictas, el orden no importa; con débiles, sí.
\item Todo ENPS es EN, no al revés; todo EBP es EBN. Inducción hacia atrás = procedimiento.
\item Gatillo: $\beta\geq(u^{\mathrm D}-u^{\mathrm C})/(u^{\mathrm D}-u^{\mathrm N})$; \emph{tit for tat} solo es ENPS con un $\beta$ concreto.
\item 2.º precio: \emph{débilmente} dominante; holandesa = 1.\textsuperscript{er} precio siempre; inglesa = 2.º solo con valores privados.
\item Deuda pública: «holandesa» = precio único y «americana» = precio ofrecido (en teoría, holandesa = descendente).
\item Reserva óptima con $U[0,1]$: $1/2$ sea cual sea $n$; supera $v_0$ (ineficiente).
\end{itemize}
\end{trampa}

\end{columnas}
\end{document}
